% called by main.tex
%

\section*{Abstract}
\addcontentsline{toc}{section}{Abstract}
\label{sec::abstract}

Anycast routing, in which the same IP address is announced from multiple locations and the Internet routing system delivers each communication to one of them, underpins a large part of today's Internet infrastructure, including Content Delivery Networks (CDNs), cloud platforms, and the third-party services embedded in mobile applications. Although anycast improves performance and resilience, it makes the geographic destination of each communication invisible to users and largely uncontrolled by service operators, since the replica that processes a request results from distributed routing decisions in which jurisdiction plays no role. This is problematic under regulations such as the General Data Protection Regulation (GDPR), which requires transparency about where personal data are processed, particularly when they are transferred to countries outside the European Economic Area (EEA). This thesis studies how anycast affects the geographic destination of communications carrying personal data, and therefore potential international transfers, and how geographic control can be restored without giving up the benefits of anycast.

First, the thesis presents Hunter, a methodology that infers the country in which an anycast communication terminates by combining traceroute analysis, last-hop identification, and latency-based multilateration. In a validation with ground truth covering 29 countries of the EEA and 2,739 measurements, Hunter achieved a precision of 97.67\% and an accuracy of 78.09\%, deliberately favouring conservative decisions that avoid attributing communications to the wrong jurisdiction.

Second, Hunter is applied to a large-scale analysis of 5,759 Android applications. Of the communications containing personal data, 80.74\% were directed towards anycast infrastructures, and they were generated by 960 applications. The results reveal a significant gap between the destinations declared and those observed: 98.65\% of the analysed applications did not adequately disclose in their privacy policies the destination countries observed, and 90\% of the Third-Party Libraries (TPLs) involved, all of which could reach destinations outside the EEA, did not disclose them either. A subsequent global campaign covering 183 countries and 504,690 routing tracks showed that this behaviour is not limited to EEA: 58.36\% of the tracks terminated in a country different from that of origin, and 99.49\% of the analysed anycast addresses exhibited cross-border routing.

Third, the thesis proposes Region-Constrained Anycast over DNS (RCA-DNS), an architecture that makes the region an explicit parameter of each communication. Each region is exposed through its own DNS name and anycast address, served only by infrastructure located within that region, so that the service selects the region according to its own policy while anycast continues to distribute traffic within it. The thesis analyses the levels at which this confinement can be enforced and shows that the architecture can be implemented on the main cloud and CDN platforms. A prototype deployed in 39 Google Cloud locations and evaluated with RIPE Atlas showed that, in well-provisioned regions such as the EEA, confinement did not increase latency: compared with a global deployment, the median round-trip time decreased by 5.99\% and the 95th percentile fell from 1,252~ms to 111~ms. In regions with little infrastructure, the latency cost is mainly explained by the lack of local locations.

Together, these contributions show that the geographic destination of anycast communications is a dynamic property that can be measured, analysed, and, through appropriate architectural mechanisms, controlled.

\newpage
\section*{Resumen}
\addcontentsline{toc}{section}{Resumen}
\label{sec::resumen}

El encaminamiento anycast, en el que una misma dirección IP se anuncia desde múltiples ubicaciones y el sistema de encaminamiento de Internet entrega cada comunicación a una de ellas, sustenta una gran parte de la infraestructura actual de Internet, incluidas las redes de distribución de contenidos (CDN), las plataformas en la nube y los servicios de terceros integrados en las aplicaciones móviles. Aunque anycast mejora el rendimiento y la resiliencia, hace que el destino geográfico de cada comunicación sea invisible para los usuarios y quede en gran medida fuera del control de los operadores de los servicios, ya que la réplica que procesa una petición resulta de decisiones de encaminamiento distribuidas en las que la jurisdicción no desempeña ningún papel. Esto resulta problemático en el marco de normativas como el Reglamento General de Protección de Datos (RGPD), que obliga a ser transparente sobre los lugares en los que se tratan los datos personales, en particular cuando se transfieren a países externos al Espacio Económico Europeo (EEE). Esta tesis estudia cómo afecta anycast al destino geográfico de las comunicaciones que contienen datos personales y, por tanto, a posibles transferencias internacionales, y cómo puede recuperarse el control geográfico sin renunciar a las ventajas de anycast.

En primer lugar, la tesis presenta Hunter, una metodología que infiere el país en el que termina una comunicación anycast combinando el análisis de traceroute, la identificación del último salto y la multilateración basada en latencias. En una validación con datos de referencia que abarcó 29 países del EEE y 2.739 medidas, Hunter alcanzó una precisión del 97,67~\% y una exactitud del 78,09~\%, priorizando deliberadamente decisiones conservadoras que evitan atribuir comunicaciones a una jurisdicción errónea.

En segundo lugar, Hunter se aplica a un análisis a gran escala de 5.759 aplicaciones Android. De las comunicaciones que contenían datos personales, el 80,74~\% se dirigían a infraestructuras anycast y fueron generadas por 960 aplicaciones. Los resultados revelan una brecha significativa entre los destinos declarados y los observados: el 98,65~\% de las aplicaciones analizadas no informaba adecuadamente en sus políticas de privacidad de los países de destino observados, y el 90~\% de las bibliotecas de terceros (TPL) implicadas, todas ellas capaces de alcanzar destinos fuera del EEE, tampoco los declaraba. Una campaña global posterior, que abarcó 183 países y 504.690 trayectorias de encaminamiento, mostró que este comportamiento no se limita al EEE: el 58,36~\% de las trayectorias terminaba en un país distinto al de origen, y el 99,49~\% de las direcciones anycast analizadas presentaba encaminamiento transfronterizo.

En tercer lugar, la tesis propone Region-Constrained Anycast over DNS (RCA-DNS), una arquitectura que convierte la región en un parámetro explícito de cada comunicación. Cada región se expone mediante su propio nombre DNS y su propia dirección anycast, servida únicamente por infraestructura situada dentro de esa región, de modo que el servicio selecciona la región según su propia política mientras anycast sigue distribuyendo el tráfico dentro de ella. La tesis analiza los niveles en los que puede aplicarse este confinamiento y muestra que la arquitectura puede implementarse en las principales plataformas en la nube y CDN. Un prototipo desplegado en 39 ubicaciones de Google Cloud y evaluado con RIPE Atlas mostró que, en regiones bien dotadas como el EEE, el confinamiento no aumentó la latencia: en comparación con un despliegue global, la mediana del tiempo de ida y vuelta disminuyó un 5,99~\% y el percentil 95 bajó de 1.252~ms a 111~ms. En regiones con poca infraestructura, el coste en latencia se explica principalmente por la falta de ubicaciones locales.

En conjunto, estas contribuciones muestran que el destino geográfico de las comunicaciones anycast es una propiedad dinámica que puede medirse, analizarse y, mediante mecanismos arquitectónicos adecuados, controlarse.
